\section{Literature Review}
In this section, we briefly survey recent work on fish classification. We first summarize the main datasets used for fish detection and classification, outlining their scope, scale, and key limitations. We then review deep learning-based methods. We will focus on the progress in convolutional neural nets, use of transfer learning, and other efforts to balance accuracy with computing efficiency. In this Fish classification divides species into three primary, distinct groups based on skeleton structure and the presence of jaws: Jawless fishes (Agnatha), Cartilaginous fishes (Chondrichthyes), and Bony fishes (Osteichthyes). Bony fish are further divided into ray-finned (Actinopterygii) and lobe-finned (Sarcopterygii) species; the classification was based on sample data from project. 

Current datasets for fish classification are the starting point for progress in the deep-sea bio identification problem. A key resource is the Fish4Knowledge dataset, which offers a large-scale basis for fish classification from composited underwater imagery but is constrained by limited environmental diversity.Kuswantori et al \cite{Kuswantori2022yolov4fish} introduced a dataset targeted at automatic sorting using an optimized YOLO-based detector, yet its narrow operational focus limits broader
ecological use. The UVOT400 dataset by Alawode et al. \cite{Alawode2023underwaterTracking} supports underwater visual tracking but remains challenged by noise and visual complexity in underwater scenes, which can impede reliable classification. The FishInTurbidWater dataset uses semi- and weakly supervised deep networks for detection in turbid videos, but sparse labels can introduce inconsistencies. Similarly, the FishNet dataset by Ma et al \cite{Ma2024semiSupervisedAquatic} advances semi-supervised species recognition for biodiversity monitoring while still facing the common constraint of insufficient high-quality labeled data. 

Early models were dominated by traditional computer vision approaches. As the field evolved, machine learning models became dominant for fish classification. Approaches using Support Vector Machines, Random Forests, and Artificial Neural Networks, consistently outperformed rule-based and purely hand-crafted feature models. This shift toward data-driven learning laid the groundwork for deep learning methods that learn discriminative representations directly from raw images. 

Deep learning fundamentally changed image classification by automating feature extraction and reducing dependence on manual descriptors. Salman et al. were among the first to apply it to unconstrained underwater fish imagery, designing a three-layer CNN to learn visual features that were then classified using SVM and kNN; despite its shallow depth, this model clearly outperformed traditional pipelines by capturing morphology and texture patterns. The introduction of AlexNet further transformed
the area by showing that deeper CNNs with ReLU, dropout, and large-scale GPU training can deliver major performance gains and inspiring applications in marine benchmarks such as LifeCLEF and SeaCLEF. Iqbal et al. \cite{Iqbal2021fishClassificationCNN} trained AlexNet from scratch for fish species classification and achieved 90.48 accuracy, while Tamou et al. \cite{Tamou2018underwaterLiveFish} used transfer learning with a pre-trained AlexNet to reach 99.45, demonstrating the strong benefits of reusing models trained on large natural image datasets for underwater environments. 

Deeper neural networks have continued to improve classification performance but introduced optimization issues, notably vanishing gradients. ResNet addressed these challenges with residual learning and identity shortcuts, enabling stable training of very deep models. While other architectures such as VGG16, InceptionV3, Xception, DenseNet, and MobileNet also perform well, ResNet has emerged as a particularly efficient and effective backbone on many image classification benchmarks. 

Available data has allowed deep learning to be the most effective solution to the fish classification problem. Researchers routinely fine-tune deep CNNs—particularly ResNet variants. For instance, Zhang et al. \cite{Zhang2022largeScaleUnderwater} introduced AdvFish, which adapts a ResNet-50 backbone with an additional loss term that suppresses background clutter and highlights key fish regions, leading to higher accuracy in high complexity scenes. 
